\documentclass[11pt]{article}
\usepackage[margin=2.6cm]{geometry}
\usepackage{mathpazo}
\usepackage{microtype}
\usepackage{booktabs}
\usepackage{graphicx}
\usepackage{float}
\usepackage{xcolor}
\usepackage{titlesec}
\usepackage{caption}
\usepackage{listings}
\usepackage[hidelinks]{hyperref}
\emergencystretch=3em

\definecolor{accent}{HTML}{1F4E79}
\definecolor{good}{HTML}{1A7F37}
\definecolor{bad}{HTML}{B42318}
\titleformat{\section}{\Large\bfseries\color{accent}}{\thesection}{0.8em}{}
\titleformat{\subsection}{\large\bfseries\color{accent!80!black}}{\thesubsection}{0.8em}{}
\captionsetup{font=small,labelfont={bf,color=accent}}
\newcommand{\armA}{\textcolor{bad}{arm A}}
\newcommand{\armB}{\textcolor{accent}{arm B}}

\lstdefinestyle{json}{
  basicstyle=\ttfamily\footnotesize,
  backgroundcolor=\color{gray!8},
  frame=single, rulecolor=\color{gray!40}, framerule=0.4pt,
  breaklines=true, columns=fullflexible, keepspaces=true,
  showstringspaces=false, aboveskip=8pt, belowskip=8pt}
\lstset{style=json}

\newsavebox{\eliBox}
\newenvironment{eli5}
  {\par\medskip\noindent\begin{lrbox}{\eliBox}\begin{minipage}{\dimexpr\linewidth-2\fboxsep-2\fboxrule}\small}
  {\end{minipage}\end{lrbox}\noindent\fcolorbox{accent!40}{accent!4}{\usebox{\eliBox}}\par\medskip}

\title{\color{accent}\bfseries Structured DAG-Diff Mutation for CARL Reasoning Chains\\
\large The Diff Language, Its Grammar Enforcement, and the 250-Mutant A/B on the Standard Pipeline (v250)}
\author{GigaEvo}
\date{2026-07-03}

\begin{document}
\maketitle

\begin{abstract}
\noindent
We replaced free-form JSON rewriting of CARL reasoning-chain genomes with a
\emph{schema-constrained positional-slot diff}: the mutation LLM emits a typed edit whose
grammar is compiled per parent set, validated with pydantic, and applied deterministically.
In a paired 250-mutant A/B on the summarizer task --- run on the framework's standard
intra-memory pipeline, with every wiring rule of the language enforced by constrained
decoding --- the diff arm persisted \textbf{zero structurally invalid programs and recorded
zero schema errors in 260 calls}; the rewrite arm persisted 8 malformed children (3.2\% of
its mutant budget), each burning a full evaluation. Fitness is at parity, and both arms now
clear their seed pools (best child 0.624 vs.\ 0.590 over seeds at 0.490/0.492). The token
trade improved over the earlier 100-mutant pair: 34\% fewer output tokens against 25\% more
input tokens. The run also surfaced the first \emph{non-grammar} failure mode: in 9 of 260
diff calls the model exhausted its 60K output budget inside the reasoning channel before
emitting the payload --- each costing one call and nothing else; nothing was persisted or
evaluated. Two earlier diagnostic runs exposed the last two grammar gaps --- a
validator-checked ordering rule, then unbounded dependency arrays --- and both closures
(narrowed per-position enums; \texttt{maxItems}) are documented as design history
(\S\ref{sec:encodings}, \S\ref{sec:depfix}). This report explains the diff language in
full, then presents the A/B evidence.
\end{abstract}

\section{Motivation}
CARL chain genomes are wire-JSON documents (\texttt{ReasoningChain.to\_dict} plus platform
extras). The stock mutation path asks the LLM to rewrite the whole document as free text:
every emission re-risks the entire syntax, and a malformed child is only discovered
\emph{after} it has been persisted and has consumed a full evaluation. The hypothesis:
constraining mutation to a typed diff language --- where structural validity is enforced by
the provider's constrained decoding and a pydantic contract --- eliminates that failure class
at the cheapest possible stage.

\section{The diff language}
\label{sec:language}
This section describes the language in its shipped form --- the fixed-key \emph{object
encoding} --- which is exactly what the A/B in \S\ref{sec:experiment} ran. How the design
arrived here (two earlier diagnostic runs, each closing one grammar gap) is told in
\S\ref{sec:encodings} and \S\ref{sec:depfix}.

\subsection{ELI5}
\begin{eli5}
Imagine the child chain as a shelf of eight numbered boxes. To describe a new chain, you
fill boxes from the left. Into each box you put either a \emph{photocopy} of a step from any
parent --- naming it by its printed label, like \texttt{a2}, optionally with corrections
scribbled on top --- or a \emph{brand-new} step you write yourself. Every box also declares
which earlier boxes it reads its input from, and here is the trick: the form for box $k$
only has checkboxes for boxes $1..k-1$. Pointing at yourself or at a later box is not
``against the rules'' --- there is literally no place on the form to write it. Box~1 has no
checkboxes at all; it always reads the raw task input. When you are done, the remaining
boxes stay empty (no gaps allowed), and whatever sits in the last filled box produces the
answer that gets graded.
\end{eli5}
\noindent
The LLM never writes a chain; it fills in this form. The form is a JSON schema compiled
fresh for each mutation from the actual parents, and the provider's constrained decoding
guarantees the reply matches the form --- a violating token can never even be sampled.

\subsection{Anatomy of a diff}
A diff \emph{is} the full child chain. One payload, real shape (parents rendered as
\texttt{a1..a2} and \texttt{b1..b3}):
\begin{lstlisting}
{
  "reasoning": "Guided Innovation; insight: faithfulness - final step.
      Insert a fact-check between draft and polish: the polish step
      paraphrases away source wording that the metric scores; checking
      the draft against the source before polishing restores it.",
  "base_parent": "A",
  "slot_1": {"kind": "keep", "id": "a1"},
  "slot_2": {"kind": "new",
             "title": "Verify facts",
             "aim": "Cross-check the draft against the source text;
                     list any unsupported claim.",
             "dependencies": ["slot_1"]},
  "slot_3": {"kind": "keep", "id": "a2",
             "edits": {"aim": "Rewrite the draft into one crisp
                       sentence, keeping verified source wording."},
             "dependencies": ["slot_2"]},
  "slot_4": null, "slot_5": null, "slot_6": null,
  "slot_7": null, "slot_8": null
}
\end{lstlisting}
Field by field:
\begin{itemize}
  \item \textbf{\texttt{reasoning}} --- a falsifiable causal hypothesis for the edit, carrying
  the evidence citations (insight anchors, intra-memory clusters, donor-step ids). Prompt
  rules forbid tautologies (``verification improves faithfulness'') and diff restatements.
  \item \textbf{\texttt{base\_parent}} --- an enum of the parent letters actually present.
  This is \emph{lineage attribution}: the parent whose run metadata and task description the
  child inherits, and the lineage credit-assignment target. It does \emph{not} restrict which
  steps may be kept.
  \item \textbf{\texttt{slot\_1} \dots\ \texttt{slot\_8}} --- the child's steps, in execution
  order. \texttt{slot\_1} is required; the rest admit \texttt{null}. Filled slots must form a
  prefix (no gaps). The last filled slot is the output step --- its answer is what gets
  scored.
\end{itemize}

\subsection{The two slot forms}
Each filled slot is one of two typed objects, discriminated by \texttt{kind}:
\begin{description}
  \item[\texttt{keep}] Reuse a rendered parent step by id. \texttt{id} is an enum compiled
  from the actual parents --- parent-prefixed (\texttt{a1}, \texttt{b3}, \dots), spanning
  \emph{all} parents, so a keep may pull a donor's step directly (crossover is free).
  Optional \texttt{edits} overlays new values onto any content field
  (\texttt{title}, \texttt{aim}, \texttt{stage\_action}, \texttt{reasoning\_questions});
  unset fields keep the parent's text. \texttt{title}/\texttt{aim} edits must be non-empty
  --- the schema carries the executor's own invariant.
  \item[\texttt{new}] A fresh LLM step: required non-empty \texttt{title} and \texttt{aim};
  optional \texttt{stage\_action} and \texttt{reasoning\_questions}.
\end{description}
Both forms take the same \texttt{dependencies} field --- except in \texttt{slot\_1}, where
the field does not exist.

\subsection{Wiring: the position-narrowed dependency enum}
\texttt{dependencies} lists the earlier slots whose outputs this step consumes --- slot
positions in the \emph{new} chain, never parent ids. The schema for slot $k$ offers the enum
$\{\texttt{slot\_1},\dots,\texttt{slot\_{k-1}}\}$ and nothing else:
\begin{center}
\footnotesize
\begin{tabular}{ll}
\toprule
slot & \texttt{dependencies} may mention \\
\midrule
\texttt{slot\_1} & --- (field absent; reads the raw task input) \\
\texttt{slot\_2} & \texttt{slot\_1} \\
\texttt{slot\_3} & \texttt{slot\_1}, \texttt{slot\_2} \\
$\vdots$ & $\vdots$ \\
\texttt{slot\_8} & \texttt{slot\_1} \dots\ \texttt{slot\_7} \\
\bottomrule
\end{tabular}
\end{center}
An empty list is legal anywhere it exists --- a step may read the raw input in parallel with
slot 1. The array is also \emph{length-capped} in the grammar (\texttt{maxItems}: $k-1$ for
slot $k$): a step can never cite more distinct predecessors than exist, and a sampler that
starts looping on a dependency token hits a forced \texttt{]} after at most seven items
instead of running away. Acyclicity of the child DAG is therefore a \emph{typing} fact, not
a checked property: every edge points strictly backwards by construction.

\subsection{Every mutation move, spelled in the language}
\begin{center}
\small
\begin{tabular}{ll}
\toprule
move & spelling \\
\midrule
edit a step & \texttt{keep} it with \texttt{edits} overriding fields \\
insert a step & put a \texttt{new} slot between two \texttt{keep}s, rewire \\
delete a step & simply never fill a slot with it \\
reorder & fill the slots in the new order \\
duplicate & \texttt{keep} the same id in two slots \\
rewire only & same \texttt{keep}s, different \texttt{dependencies} \\
crossover & \texttt{keep} ids from more than one parent \\
full rewrite & every slot \texttt{new} \\
branching / parallel DAG & two slots share a dependency; a later slot joins them \\
\bottomrule
\end{tabular}
\end{center}
Two of these in miniature. Deletion is pure omission --- collapsing a 2-step parent to its
final step is a one-slot diff:
\begin{lstlisting}
{"reasoning": "...", "base_parent": "A",
 "slot_1": {"kind": "keep", "id": "a2"},
 "slot_2": null, ..., "slot_8": null}
\end{lstlisting}
Crossover is just a keep-id from another parent --- graft B's fact-extractor in front of A's
polisher:
\begin{lstlisting}
{"reasoning": "... parent:B step:b1 ...", "base_parent": "A",
 "slot_1": {"kind": "keep", "id": "b1"},
 "slot_2": {"kind": "keep", "id": "a2", "dependencies": ["slot_1"]},
 "slot_3": null, ..., "slot_8": null}
\end{lstlisting}

\subsection{What the grammar makes impossible}
The design goal: a defect class should be \emph{unrepresentable}, not detected. What an
unconstrained rewriter routinely gets wrong, this language cannot even spell:
\begin{table}[H]\centering\small
\begin{tabular}{ll}
\toprule
illegal child & why it cannot be emitted \\
\midrule
self-dependency (\texttt{slot\_k} $\to$ \texttt{slot\_k}) & slot $k$'s enum stops at \texttt{slot\_{k-1}} \\
forward dependency (\texttt{slot\_2} $\to$ \texttt{slot\_5}) & same enum narrowing \\
any dependency on slot 1 & the field does not exist there; extra keys rejected \\
dependency cycle & impossible transitively: all edges point backwards \\
keep of a nonexistent step (\texttt{a9}) & \texttt{id} enum is compiled from the actual parents \\
9-step chain & there is no \texttt{slot\_9} property \\
empty chain & \texttt{slot\_1} is required \\
empty \texttt{title}/\texttt{aim} (new or edit) & \texttt{minLength: 1} in the schema \\
runaway \texttt{dependencies} list (\texttt{slot\_1} spammed 70$\times$) & \texttt{maxItems}: $k-1$ caps the array \\
invented fields, malformed JSON & \texttt{additionalProperties: false}; constrained decoding \\
\midrule
gap fill (\texttt{slot\_3} filled, \texttt{slot\_2} null) & \emph{the} residual pydantic validator (contiguity) \\
\bottomrule
\end{tabular}
\caption{Everything above the line is enforced by the decoding grammar itself. Contiguity is
the single rule left to a validator --- and combined with the enums it also makes
``dependency on an empty slot'' impossible.}
\end{table}

\subsection{How enforcement works}
\begin{enumerate}
  \item \textbf{Per-mutation compilation.} \texttt{AllowedDagChanges.build\_schema} parses
  the actual parents and generates a pydantic model with \texttt{create\_model}: the keep-id
  enum, the \texttt{base\_parent} enum, and eight per-position slot unions are all derived
  from what really exists right now.
  \item \textbf{Constrained decoding.} The JSON schema ships to the provider as the
  structured-output contract (\texttt{method="json\_schema"}). The provider compiles it into
  a decoding grammar; tokens that would violate the schema get zero probability. The model
  cannot emit a forward reference for the same reason it cannot emit an unbalanced brace.
  \item \textbf{Portable subset.} Provider grammar compilers accept different JSON-schema
  dialects; Gemini is the strictest we probed (hard-rejects \texttt{\$ref}/\texttt{\$defs}
  and \texttt{const}). \path{gigaevo/llm/schema_compat.py} emits every schema in the
  lowest-common-denominator subset: refs inlined, \texttt{const} $\to$ single-value
  \texttt{enum}, decorative keys dropped. Nothing is provider-conditional at runtime.
  \item \textbf{Grammar budget.} Strict compilers limit construct complexity, not bytes (a
  padded 120KB schema passes; an 87KB slot-union grammar does not). Branching the slot
  models per parent blew the cliff at 2 parents $\times$ 8 slots ($\sim$27KB); the shipped
  encoding uses a \emph{single} model --- one global parent-namespaced keep-id enum,
  \texttt{base\_parent} as a plain enum field, no union discriminator --- at 13.7KB with
  $\sim$2$\times$ headroom.
  \item \textbf{Pydantic backstop.} The same model validates the reply before anything
  touches the engine; the contiguity rule lives here. A rejection costs one LLM call ---
  nothing is persisted, nothing is evaluated.
\end{enumerate}

\subsection{From diff to child}
The applier is deterministic --- the LLM's creative surface ends at the diff:
transcription walks \texttt{slot\_1..slot\_8} until the first \texttt{null}; \texttt{keep}
slots copy content fields from the referenced parent step (any parent) and overlay
\texttt{edits}; \texttt{dependencies} become sorted, deduplicated integer positions; run
metadata and the task description come from \texttt{base\_parent}; the last step is stamped
as the output step. The result is round-tripped through \texttt{ReasoningChain} as a final
assertion and emitted as canonical wire JSON --- the child is indistinguishable in format
from a hand-written genome.

\section{Implementation}

\subsection{Operator and vocabulary seam}
Two pieces, both behind existing seams:
\begin{itemize}
  \item \path{gigaevo/evolution/mutation/structured_diff.py} --- the generic engine-level
  operator \texttt{StructuredDiffMutationOperator}: it asks an \texttt{AllowedChanges}
  vocabulary to build a \texttt{DiffSchema} for the current parents, runs a
  structured-output LLM call against that schema, validates, and applies. It knows nothing
  about chains.
  \item \path{gigaevo/chains/dag_changes.py} --- \texttt{AllowedDagChanges}, the
  chain-specific vocabulary of \S\ref{sec:language}: schema compilation, parent rendering
  (the \texttt{a1}/\texttt{b2} listings the LLM sees), the applier, and the natural-language
  contract injected into the mutation prompt.
\end{itemize}

\subsection{The encoding search}
\label{sec:encodings}
Getting the diff grammar through the strictest compiler took four encodings:
\begin{table}[H]\centering\footnotesize
\begin{tabular}{llll}
\toprule
encoding & ordering rule in grammar? & size & Gemini verdict \\
\midrule
E1: tuple-union (\texttt{prefixItems}) & yes (per-slot models) & $\sim$87KB & rejected (complexity) \\
E2: cons-list recursion (\texttt{\$ref}) & yes (per-depth) & --- & rejected (\texttt{\$ref} banned) \\
E3: flat \texttt{anyOf[keep,new]} array & \textcolor{bad}{no --- pydantic validator} & $\sim$8KB & accepted --- superseded \\
E4: fixed-key object (\texttt{slot\_k}) & \textcolor{good}{yes (narrowed enums + \texttt{maxItems})} & 13.7KB & \textcolor{good}{accepted --- \textbf{shipped, ran the A/B}} \\
\bottomrule
\end{tabular}
\caption{Raw size is not the limit; construct complexity is. E4 is the encoding described in
\S\ref{sec:language} and the one the A/B ran; E3 differed only in expressing slots as an
array and checking the ordering rule post-decode (\S\ref{sec:depfix}).}
\end{table}

\subsection{Engine integration}
The engine assumed the mutator's \texttt{base\_parent} is a 1-based integer (wire-JSON
convention); diff payloads use namespace letters, so every diff mutant died pre-persist in
the first smoke (\texttt{int('A')}). Fixed test-first with \texttt{base\_parent\_index()} in
\texttt{gigaevo/evolution/engine/mutation.py} (int or namespace letter; garbage falls back to
parent 1 with a warning). Other pre-launch fixes (config-group promotion, loader pattern
clobber, prompt-fairness rewrite for arm A, token accounting) are in
\texttt{04\_issues\_log.md}.

\section{Experiment}
\label{sec:experiment}
Paired 250-mutant runs, identical seeds (4), engine config, and models: mutation
gemini-3.5-flash (OpenRouter, reasoning effort high, 60K max output tokens); chain executor
Qwen3-235B-Instruct (LiteLLM proxy). Task: \texttt{problems/chains/summarizer}, ROUGE-L
fitness. \textbf{\armA{}} = stock full wire-JSON rewrite with a chain-faithful prompt;
\textbf{\armB{}} = structured slot diff on the shipped object encoding E4 of
\S\ref{sec:language}, \texttt{maxItems} included.

\paragraph{Pipeline.} Both arms run the framework's \emph{standard intra-memory pipeline}:
every evaluated parent carries a descriptive lineage card (\texttt{IntraMemoryStage}) and
prescriptive insights (\texttt{MutationSuggestionStage}), and both arms' mutation prompts
receive the identical per-parent evaluation-context block, so the only cross-arm difference
is the mutation operator. All earlier runs in this campaign --- including the 100-mutant
pair (\S\ref{sec:full100}) --- had run on the legacy pipeline base;
\texttt{JsonChainPipelineBuilder} was rebased onto the standard builder immediately before
this pair, so this report measures the operator under the framework's current default
context harness.

\section{Results}

\subsection{Validity funnel (headline)}
\begin{table}[H]\centering\small
\begin{tabular}{lcc}
\toprule
 & \armA{} (rewrite) & \armB{} (diff) \\
\midrule
mutation LLM calls & 256 & 260 \\
cancelled in flight at shutdown (0 tokens) & 2 & 0 \\
generation failures (pre-persist) & \textcolor{bad}{4} (no code extracted) & \textcolor{bad}{10} (9 length + 1 transient) \\
\quad\emph{of which schema errors} & --- & \textbf{\textcolor{good}{0}} \\
children persisted & 250 & 250 \\
evaluation cut by the shutdown sweep & 0 & 2 \\
\textbf{invalid after evaluation} & \textbf{\textcolor{bad}{8}} (7 json parse, 1 undecodable) & \textbf{\textcolor{good}{0}} \\
valid children & 242 & 248 \\
valid yield per completed evaluation & 96.8\% & \textbf{100\%} \\
valid yield per LLM call & 94.5\% & 95.4\% \\
\bottomrule
\end{tabular}
\caption{The arms fail at different stages with very different costs. The ``shutdown''
rows are artifacts of the fixed mutant budget, not failures: when the budget fills,
in-flight mutation calls and evaluations are swept. (Separately from the operator, the
pipeline's suggestion stage had 7/\armA{} and 8/\armB{} in-flight calls swept at shutdown;
they touch neither this funnel nor the token table.)}
\end{table}
\noindent
\armA{} defects pass mutation silently (the operator extracts a string) and die in
evaluation: 8 mutants --- 3.2\% of the budget --- persisted malformed JSON (7 parse
errors, 1 undecodable output), each burning a mutant slot plus a full evaluation pipeline.
\armB{} persisted \textbf{nothing} invalid: 260 calls, zero schema rejections, zero invalid
children --- the grammar guarantee of the 100-mutant pair (\S\ref{sec:full100}) held
unchanged at 2.5$\times$ scale. What \armB{} did pay is a new, non-grammar failure class
(\S\ref{sec:lengthfail}): 9 calls exhausted the 60K output budget inside the reasoning
channel and 1 returned a transiently unparsable response --- each cost one LLM call,
nothing was persisted, nothing was evaluated. The failure asymmetry is the point:
rewrite-arm defects cost an \emph{evaluation}; diff-arm defects cost a \emph{call}.
The 2 \armB{} children whose evaluations were shutdown-cut both parse as valid chains with
strictly-backward wiring; they are excluded from the yield denominators.

\subsection{The residual failure mode: reasoning-budget exhaustion}
\label{sec:lengthfail}
Nine of \armB's 260 calls (3.5\%) failed with the SDK's
\texttt{LengthFinishReasonError}: gemini-3.5-flash, at reasoning effort high, burned the
entire 60K max-token budget in its private reasoning channel before emitting the first
payload token, so constrained decoding never produced a document. The failures were spread
through the run (not an endsweep artifact) and cost $\sim$4 minutes of latency each. This
is not a grammar gap --- the schema constrains \emph{what} the model emits, not \emph{how
long it thinks} --- and it plausibly shares a root cause with \armA's truncation-flavored
parse errors: when generation hits the length ceiling, structured output fails the call
cleanly at zero downstream cost, while free-form rewriting hands a truncated document to
the engine, which persists and evaluates it. Mitigations are policy, not language: raise
the output budget, cap reasoning effort, or retry on length.

\subsection{Fitness --- parity, with both arms clearing their seeds}
Child mean 0.400 (\armA) vs.\ 0.421 (\armB); medians 0.407 vs.\ 0.428. Unlike the
100-mutant pair --- where \armB's best child lost to its own seed --- both arms now beat
their seed pools: best child 0.624 (\armA) and 0.590 (\armB) against seed bests of
0.490/0.492 (the same 4 seeds, re-scored per run; the gap between the two seed numbers is
pure executor sampling noise). \armB's mean/median edge ($\sim$+0.02) and \armA's
best-child edge ($\sim$+0.03) both sit inside single-replicate noise --- no same-config
variance floor exists for this problem yet --- so we claim parity with healthy search in
both arms, nothing more.
\begin{figure}[H]\centering
\includegraphics[width=\textwidth]{ab_overview.png}
\caption{Left: per-child fitness (dots) and best-so-far (line). Right: validity funnel.}
\end{figure}

\subsection{Structural exploration}
\armA{} spans 1--4 steps ($\{46, 126, 55, 15\}$ across lengths 1/2/3/4); \armB{} spans
1--3 ($\{29, 162, 57\}$), concentrating hardest on 2-step chains (65\% of its valid
children). At this scale the \emph{rewrite} arm explored marginally longer chains; the
language admits up to 8 slots but neither arm exercises the headroom, so we claim no
structural-exploration advantage in either direction --- the grammar's guarantee is only
that whatever lengths arrive are well-formed.
\begin{figure}[H]\centering
\includegraphics[width=0.7\textwidth]{ab_nsteps.png}
\caption{Chain-length distribution of valid children.}
\end{figure}

\subsection{Cost}
\begin{table}[H]\centering
\begin{tabular}{lcc}
\toprule
mutation side (gemini-3.5-flash) & \armA{} & \armB{} \\
\midrule
tokens in & 2{,}091{,}237 & 2{,}609{,}025 ($+25\%$) \\
tokens out & 2{,}218{,}918 & \textbf{1{,}471{,}961 ($-34\%$)} \\
mean output tokens per completed call & 8{,}876 & 5{,}888 ($-34\%$) \\
\bottomrule
\end{tabular}
\caption{The trade improved under the standard pipeline: $-34\%$ output (vs.\ $-11\%$ in
the 100-mutant pair) against $+25\%$ input (vs.\ $+46\%$) --- the fixed $\sim$3.4K-token
schema contract is a smaller fraction of the now-larger per-call prompt, which carries the
lineage card and suggester insights in both arms. Output includes the model's reasoning
channel (effort high), which varies run to run. Caveat: the table counts recorded usage
only; the 9 length-exhausted \armB{} calls (and \armA's 6 zero-usage calls) recorded zero
tokens but billed unaudited server-side reasoning --- if each \armB{} failure billed the
full 60K cap, the output saving compresses to roughly $-9\%$. Executor completion tokens:
235K (\armA) vs.\ 230K (\armB).}
\end{table}
\begin{figure}[H]\centering
\includegraphics[width=\textwidth]{ab_tokens.png}
\caption{Per-mutant token spend. Left: output tokens per mutation call (dots) with a 10-call
rolling mean (line) --- \armA{} hovers at 8--10K tokens per call, \armB{} at 4--6K with a
few reasoning-channel spikes up to 60K. Zero-token points are calls that returned no
completion (\armA's extraction failures, \armB's length failures, or in-flight calls
cancelled when the mutant budget was reached). Right: cumulative mutation tokens; solid
lines show the $-34\%$ output gap, dashed lines \armB's input premium from shipping the
schema with every call.}
\end{figure}

\section{Design history: the runs that hardened the grammar}
\label{sec:depfix}
The grammar the A/B ran was hardened by two earlier 100-mutant diagnostic runs, each of
which exposed exactly one failure class and forced it into the schema, followed by a
100-mutant confirmation pair on the shipped grammar. The requirement was explicit
throughout: fix defects in the grammar, not post-hoc.

\subsection{Gap 1: the ordering rule lived in a validator (encoding E3)}
The first run used E3 --- slots as a flat array, the \emph{dependencies-point-backwards}
rule checked by pydantic after decoding. All three of its pre-persist rejections (3/108
calls) were the same degenerate violation: slot~1 emitting
\texttt{dependencies:\ ["slot\_1"]}, a self-reference on the first slot --- precisely the
one rule the grammar did not carry. The fix is the fixed-key object encoding of
\S\ref{sec:language}: \texttt{slot\_1..slot\_8} as separate properties, each slot's
dependency enum narrowed to earlier slots, slot~1 with no dependencies field at all.
Self- and forward-references became unrepresentable; the sole surviving validator is
contiguous fill. Probes shaped the final form: Gemini rejects the per-parent-branched
variant (grammar cliff between 7 and 8 slots) and nested union discriminators, but accepts
the single-model variant --- keep-ids in one global parent-namespaced enum,
\texttt{base\_parent} a plain enum --- at 13.7KB with $\sim$2$\times$ headroom. A welcome
side effect: keeps from \emph{any} parent are now first-class, so crossover no longer
requires re-creating donor steps. Design notes:
\texttt{experiments/carl\_dag\_diff\_dependency\_rule\_fix.md}.

\subsection{Gap 2: unbounded arrays outrun the provider's enum enforcement}
The second run, on the object encoding, produced zero ordering violations --- and two
rejections of a class nobody had predicted: \emph{degenerate repetition}. The sampler
locked into a loop and emitted \texttt{"slot\_1"} $\sim$70 times inside one
\texttt{dependencies} array; deep into the list, items came back as \texttt{"slot"} and
\texttt{""} --- the provider's enum enforcement demonstrably lapses far into a long array.
The pydantic backstop caught both replies (nothing was persisted), but the calls were
wasted. The grammar fix: \texttt{maxItems}: $k-1$ on slot $k$'s dependency array
(\texttt{Field(max\_length=k-1)}). A repetition loop now hits a forced \texttt{]} after at
most seven items --- well inside the enforcement horizon --- and a step can never cite more
distinct predecessors than exist. Everything after this fix ran with the cap: the
confirmation pair below and the 250-mutant A/B of \S\ref{sec:experiment} recorded
\textbf{zero} schema errors in 106 and 260 calls respectively.

\subsection{The 100-mutant confirmation pair, and the pipeline rebase}
\label{sec:full100}
The shipped grammar's first paired test (\texttt{full100}, 2026-07-03) ran 100 mutants per
arm: the diff arm completed with zero failures of \emph{every} kind --- no schema
rejections, no call errors, nothing invalid persisted (106 calls) --- while the rewrite arm
persisted 5 malformed children and lost 2 calls to failed extraction; fitness sat at
parity, and the token bill measured $-11\%$ output / $+46\%$ input. That pair, like every
run before it, used the problem's original pipeline builder, which subclassed the
\emph{legacy} base (flat insights and lineage-summary stages). Before the 250-mutant pair,
\texttt{JsonChainPipelineBuilder} was rebased onto the standard intra-memory builder ---
lineage cards plus prescriptive mutation suggestions, the framework's current default ---
with the two JSON-parser stage swaps as its only remaining delta; a 5-mutant smoke
confirmed both new stages fire with zero errors. The v250 pair of \S\ref{sec:experiment}
is therefore the shipping-configuration measurement; \texttt{full100} stands as the
grammar-confirmation run.

\subsection{Verification}
\begin{itemize}
  \item \textbf{Unit:} 32 tests in \path{tests/chains/test_dag_changes.py} --- every mutation
  archetype constructs a valid child; eleven illegal-payload classes (self-dep, forward dep,
  slot-1 dep, dangling id, gap fill, empty chain, empty aim, 9th slot, over-cap dependency
  list, the 71-item degenerate repetition payload itself, \dots) are rejected; schema-shape
  tests walk the emitted JSON schema and assert the per-position enum narrowing and the
  \texttt{maxItems} caps; a 300-case fuzz applies random schema-valid diffs and round-trips
  every child.
  \item \textbf{Suite:} targeted run green (chains + integration, 203 passed); the full
  1740-test sweep passed on the object encoding; lint clean.
  \item \textbf{Live smoke} (5 mutants, gemini-3.5-flash, same launch path as the A/B):
  5/5 mutation calls grammar-accepted, zero schema rejections, all children valid ---
  confirming Gemini accepts \texttt{maxItems} before the full run.
\end{itemize}

\section{Verdict}
The hypothesis is \textbf{supported at 250-mutant scale on the standard pipeline}: invalid
persisted programs $8\to0$, valid yield per completed evaluation $96.8\%\to100\%$, zero
schema errors in 260 calls, fitness at parity with both arms clearing their seeds. The
token trade improved to $-34\%$ output against $+25\%$ input (with the server-side billing
caveat on failed calls). Every wiring rule of the diff language is unrepresentable-to-break
--- 366 consecutive constrained calls across the confirmation pair and this A/B without a
single schema error --- and the only validator left is contiguity. The residual failure
mode is no longer structural: reasoning-budget exhaustion on 3.5\% of calls
(\S\ref{sec:lengthfail}), addressable by output-budget, effort, or retry policy rather than
by grammar. Next: adopt E4 as the standing encoding for chain problems, add a
retry-on-length policy, and let longer runs test whether the 8-slot headroom converts to
structural exploration and fitness.
\end{document}
